\documentclass[11pt,a4paper]{article}
\usepackage[margin=29mm]{geometry}
\usepackage[T1]{fontenc}
\usepackage{lmodern}
\pdfgentounicode=1
\pdfglyphtounicode{mapsto}{007C}
\pdfglyphtounicode{parenleftbig}{0028}
\pdfglyphtounicode{parenrightbig}{0029}
\pdfglyphtounicode{parenleftBigg}{0028}
\pdfglyphtounicode{parenrightBigg}{0029}
\pdfglyphtounicode{bracketleftBigg}{005B}
\pdfglyphtounicode{bracketrightBigg}{005D}
\pdfglyphtounicode{circleplusdisplay}{2A01}
\pdfglyphtounicode{summationtext}{2211}
\pdfglyphtounicode{producttext}{220F}
\pdfglyphtounicode{intersectiontext}{22C2}
\pdfglyphtounicode{summationdisplay}{2211}
\pdfglyphtounicode{productdisplay}{220F}
\pdfglyphtounicode{radicalbigg}{221A}
\usepackage{amsmath,amssymb,amsthm,mathtools}
\usepackage{enumitem}
\usepackage{needspace}
\usepackage{tikz}
\usetikzlibrary{arrows.meta,fit,backgrounds}
\usepackage{microtype}
\usepackage{xcolor}
\usepackage[colorlinks=true,linkcolor=blue!50!black,citecolor=blue!50!black,urlcolor=blue!50!black]{hyperref}
\usepackage[all]{hypcap}
\hypersetup{pdftitle={A counterexample to Kaplansky's quasitrace conjecture and failure of tensor-product stable finiteness},pdfauthor={OpenAI}}
\newtheorem{theorem}{Theorem}[section]
\newtheorem{lemma}[theorem]{Lemma}
\newtheorem{proposition}[theorem]{Proposition}
\newtheorem{corollary}[theorem]{Corollary}
\theoremstyle{definition}
\newtheorem{definition}[theorem]{Definition}
\theoremstyle{remark}
\newtheorem{remark}[theorem]{Remark}
\numberwithin{equation}{section}
\setlength{\emergencystretch}{1.5em}
\title{A counterexample to Kaplansky's quasitrace conjecture and failure of tensor-product stable finiteness}
\author{OpenAI}
\date{September 23, 2026}
\begin{document}
\maketitle
\begin{abstract}
We refute Kaplansky's quasitrace conjecture by constructing a separable
unital complex \(C^*\)-algebra that admits normalized \(2\)-quasitraces
but has no tracial state. As a consequence, we show that the minimal tensor
product of two unital simple stably finite complex \(C^*\)-algebras can be
properly infinite, even when one factor is the reduced free-group algebra
\(C_r^*(\mathbb F_2)\).
\end{abstract}

\section{Introduction}
\label{sec:introduction}

A unital \(C^*\)-algebra is \emph{stably finite} when every matrix
algebra over it has no nonunitary isometry, and is \emph{properly
infinite} when it contains two isometries with orthogonal ranges.
Milh{\o}j and R{\o}rdam asked whether the minimal tensor product of two
unital simple stably finite \(C^*\)-algebras must again be stably finite
\cite[Question~2.5]{MR}. They proved that this question is equivalent
to the quasitrace problem and that a negative answer to that problem would yield a
properly infinite tensor product with \(C_r^*(\mathbb F_\infty)\)
\cite[Section~2 and Theorem~3.1]{MR}. We construct the required
quasitrace counterexample. The resulting tensor product can be properly
infinite with one factor equal to the reduced \(C^*\)-algebra
\(C_r^*(\mathbb F_2)\) of the free group on two generators.

A quasitrace retains positivity and the identity
\(\tau(x^*x)=\tau(xx^*)\), but is required to be linear only on
commutative subalgebras, together with compatibility with real and
imaginary parts. A \(2\)-quasitrace must also extend to the two-by-two
matrix algebra. Whether this extra condition forces full linearity
is the general quasitrace problem. It is equivalent to Kaplansky's
question whether every \(AW^*\)-factor of type \(\mathrm{II}_1\) is a
von Neumann algebra; see \cite[Section~2]{MR} and
\cite[Conjectures~1.1--1.2 and Theorem~4.11]{Gow}.
Here the term \emph{von Neumann algebra} has its abstract \(W^*\)
meaning: the \(C^*\)-algebra admits a Banach-space predual.

Gow's historical account explains that Kaplansky introduced
\(AW^*\)-algebras in 1951 to characterize \(W^*\)-algebras by intrinsic
algebraic conditions, independent of a chosen Hilbert-space
representation. The factor types have different outcomes: every type
\(\mathrm{I}\) \(AW^*\)-factor is a \(W^*\)-algebra, whereas type
\(\mathrm{III}\) counterexamples exist. The type \(\mathrm{II}_1\)
factor question is the finite case addressed here
\cite[Introduction]{Gow}.

Quasitraces also connect the projection-theoretic origins of the question
to dimension theory for general \(C^*\)-algebras. Cuntz developed
dimension functions for simple algebras \cite{Cuntz1978}, and
Blackadar and Handelman established the correspondence between
\(2\)-quasitraces and lower semicontinuous dimension functions
\cite{BH}. Haagerup credits Blackadar and Handelman with proving that
a positive answer to Kaplansky's factor question would imply that every
\(2\)-quasitrace on a unital \(C^*\)-algebra is a trace. He also recalls
their realization of such quasitraces as pullbacks of quasitraces on
finite \(AW^*\)-algebras along unital \(*\)-homomorphisms
\cite[Introduction and Proposition~3.3]{Haagerup}.
Blackadar and R{\o}rdam later developed the extension of states on
preordered semigroups that underlies the general existence criterion
\cite{BR}.

The exact case was settled by Haagerup: every \(2\)-quasitrace on a
unital exact \(C^*\)-algebra is a trace \cite[Theorem~5.11]{Haagerup}.
A first draft was completed in 1991 and circulated thereafter; the
paper was published in 2014. Without
exactness, the existence of quasitraces is governed by a different
condition. The Cuntz--Blackadar--Handelman existence theorem provides a
normalized \(2\)-quasitrace whenever no matrix algebra over the given
unital algebra is properly infinite \cite[Theorem~2.3]{MR}.
Thus a unital algebra with that property but without tracial states
would answer the general question negatively.

We construct such an algebra and obtain a fixed positive pair
witnessing nonadditivity. Throughout, a normalized \(2\)-quasitrace
has value one on the unit of its domain; its corner-preserving
extension to \(M_2\) has value two on the matrix identity.
The definitions are recalled in Section~\ref{sec:preliminaries}.

\Needspace{10\baselineskip}
\begin{theorem}
\label{thm:main}
There exist a separable unital complex \(C^*\)-algebra \(A\), admitting
a normalized \(2\)-quasitrace, and positive contractions \(a,b\in A\)
such that every normalized \(2\)-quasitrace \(\tau\) on \(A\) satisfies
\[
 \tau(a+b)-\tau(a)-\tau(b)\ge\frac1{144}.
\]
\end{theorem}

Theorem~\ref{thm:main} refutes the conjecture that every normalized
\(2\)-quasitrace on every unital complex \(C^*\)-algebra is a tracial
state. By the equivalence recalled above, it also gives a negative
answer to Kaplansky's \(\mathrm{II}_1\) \(AW^*\)-factor question.
The constructed algebra is necessarily nonexact, by Haagerup's
result.

The matrix-extension requirement is essential to this conclusion.
Haagerup reports Kirchberg's 2006 unpublished example of a
\(1\)-quasitrace that is not a \(2\)-quasitrace, and hence is not a trace
\cite[Introduction]{Haagerup}. This example lacks the matrix extension
required here.
Gow's recent work places the established factor--quasitrace equivalence
alongside further equivalent formulations involving maximal abelian
subalgebras and bidual retractions \cite[Theorem~4.11]{Gow}.

The tensor-product conclusion and two related consequences can be stated
together. All tensor products below are minimal \(C^*\)-tensor products.

\Needspace{12\baselineskip}
\begin{corollary}[Tensor products and stable finiteness]
\label{cor:stable-finiteness}
\leavevmode
\begin{enumerate}[label=\textup{(\roman*)}]
\item There are two unital simple stably finite complex \(C^*\)-algebras
whose minimal tensor product is properly infinite. One factor can be
taken to be \(C_r^*(\mathbb F_2)\).
\item There is a separable unital complex \(C^*\)-algebra \(B\) admitting
a normalized \(2\)-quasitrace such that
\(B\otimes_{\min}C_r^*(\mathbb F_2)\) is properly infinite and admits no
normalized \(2\)-quasitrace.
\item There is a separable unital stably finite complex \(C^*\)-algebra
with no tracial state.
\end{enumerate}
\end{corollary}

The three assertions use related but different algebras. The simple
tensor-product example comes from a non-von-Neumann \(AW^*\)-factor of
type \(\mathrm{II}_1\), with no separability assertion for that factor.
The quasitrace-loss example uses the explicit separable algebra \(B\)
constructed below; we prove that no matrix algebra over \(B\) is properly
infinite, not that \(B\) is stably finite. A finitely generated subalgebra
of the \(AW^*\)-factor gives the separable stably finite traceless example,
without a simplicity assertion. The short deductions from established
trace-obstruction and tensor-product results are given in
Section~\ref{sec:consequences}.

\subsection*{The construction}

We first build a unital algebra \(B=C^*(1,x_1,\ldots,x_{24})\) such that
\begin{equation}
 \sum_{i=1}^{24}x_i^*x_i=1,\qquad
 \sum_{i=1}^{24}x_ix_i^*\le\tfrac23\,1,
 \label{eq:intro-relations}
\end{equation}
while no \(M_m(B)\), \(m\ge1\), contains two isometries with orthogonal
ranges. A tracial state would make the two sums in
\eqref{eq:intro-relations} have the same value, so \(B\) has none.
These relations instantiate Haagerup's characterization of the absence
of tracial states by a finite tuple with column sum one and a strict
row bound \cite[Lemma~2.1]{Haagerup}.
Milh{\o}j and R{\o}rdam proposed separating a uniformly bounded length
of such a tuple from increasingly delayed matrix proper infiniteness
\cite[Remark~4.7]{MR}. The construction below uses a bounded product
and a generated subalgebra, with a separate test for every matrix size
and polynomial degree. The fixed number of generators preserves the
trace obstruction while all finite polynomial candidates are excluded.

The generators are shifts between finite-rank vector bundles. For
each pair of positive integers \((m,R)\), we choose a compact base
\(Y_{m,R}\) and bundles \(E_j\to Y_{m,R}\) for all \(j\geq0\), with
\(E_0\) trivial. Only levels \(0,\ldots,R\) enter the finite topological
test. A polynomial matrix \(P\in M_{m,2m}(B)\) whose entries contain
only words of length at most \(R\) restricts to a bundle map
\[
 P_0:E_0^{\oplus2m}\longrightarrow
 \left(\bigoplus_{j=0}^{R}E_j\right)^{\oplus m}.
\]
The target's Chern class forbids the trivial image bundle that would
result from \(P^*P\geq\tfrac12 I\). Write \(S_i^{m,R}(y)\) for the
corresponding forward shift on the Hilbert direct sum
\(\bigoplus_{j\geq0}E_j(y)\). Uniform norm bounds allow the simultaneous definitions
\[
 x_i=\bigl(S_i^{m,R}(y)\bigr)_{m,R\geq1,\;y\in Y_{m,R}},
 \qquad B=C^*(1,x_1,\ldots,x_{24}),
\]
in one bounded operator product. The number of generators and the row
bound in \eqref{eq:intro-relations} are uniform, while the bundle ranks
and base dimensions may depend on \((m,R)\). A polynomial matrix
approximating a candidate pair of isometries has the same coefficients
at every coordinate. Its finite degree selects one of the tests already
present in the product, which obstructs that approximation.

The topological obstruction follows the use of characteristic classes
in the constructions of Villadsen \cite[Section~3]{Villadsen1999}
and R{\o}rdam \cite[Proposition~4.5]{Rordam2003}.
Here projective-space factors retain a nonzero Chern class after a
controlled trivial summand is allowed. The parameters are selected
anew for each finite test.

Two ingredients keep the expansion and the topological obstruction
compatible. First, a family of bundle maps whose collective image
expands every source subspace by a factor of two can be normalized
continuously to have column sum one and row sum at most \(2/3\).
The determinant-potential approach is related to operator scaling
\cite[Section~4]{Gurvits2003} and to geodesically convex matrix
optimization \cite[Corollary~2.11]{SraHosseini2015}.
We prove strict convexity and existence of a unique minimizer for the
present rectangular maps; uniqueness makes the normalization compatible
with changes of bundle frame.
Second, the bundle tests separate small and large source subspaces:
sixteen maps into reduced copies of old line bundles handle the
small ones, while eight generic maps into a new line-bundle block
handle the large ones. This keeps the generator count fixed as the
matrix size and word-length bound increase.

Finally, set \(A=M_{48}(B)\) and apply Theorem~\ref{thm:existence}
directly to \(A\). A finite telescoping construction, using only
linearity on commuting elements, gives the pair in
Theorem~\ref{thm:main}. This last step works for every normalized
\(1\)-quasitrace on \(A\), hence for every normalized \(2\)-quasitrace
there, and requires no formula for a chosen quasitrace.

\subsection*{Organization}

Section~\ref{sec:preliminaries} records the operator-algebra
conventions and the existence theorem. Section~\ref{sec:bundle-map-lemmas}
proves the incidence and scaling lemmas. The finite bundle tests
are assembled in Section~\ref{sec:bundle-tests}. Section~\ref{sec:algebra}
constructs \(B\), excludes proper infiniteness in every matrix size,
exhibits the positive pair, and completes the proof of
Theorem~\ref{thm:main}. Section~\ref{sec:consequences} proves
Corollary~\ref{cor:stable-finiteness}.

\section{Quasitraces and proper infiniteness}
\label{sec:preliminaries}

All \(C^*\)-algebras in this paper are complex, and all matrix
algebras carry their usual \(C^*\)-norms. For a unital \(C^*\)-algebra
\(C\), a \emph{\(1\)-quasitrace} is a function
\(\tau:C\to\mathbb C\) satisfying
\begin{enumerate}[label=\textup{(\roman*)}]
\item \(\tau(x^*x)=\tau(xx^*)\ge0\) for every \(x\in C\);
\item \(\tau(h+ik)=\tau(h)+i\tau(k)\) for self-adjoint \(h,k\in C\);
\item the restriction of \(\tau\) to every abelian \(C^*\)-subalgebra
of \(C\) is complex linear.
\end{enumerate}
It is a \emph{\(2\)-quasitrace} if a \(1\)-quasitrace on \(M_2(C)\)
restricts to \(\tau\) on the upper-left corner. It is
\emph{normalized} when \(\tau(1_C)=1\).
The extension then has value two at \(1_{M_2(C)}\): the two
diagonal corner units have the same value by \textup{(i)} and
their values add by \textup{(iii)}.

A unital \(C^*\)-algebra \(C\) is \emph{properly infinite} when it
contains isometries \(s_1,s_2\) with orthogonal ranges. Equivalently,
\[
 s_i^*s_j=\delta_{ij}1_C,\qquad i,j\in\{1,2\}.
\]
It is \emph{stably properly infinite} if \(M_m(C)\) is properly
infinite for some positive integer \(m\). Absence of stable proper
infiniteness is the condition needed below; no assertion about
finiteness of every projection is required.

\begin{theorem}[Quasitrace existence]
\label{thm:existence}
A unital \(C^*\)-algebra \(C\) admits a normalized \(2\)-quasitrace
if and only if \(M_m(C)\) is not properly infinite for every
positive integer \(m\).
\end{theorem}

We use Theorem~\ref{thm:existence} in the formulation of
\cite[Theorem~2.3]{MR}, with the normalization and matrix-extension
conventions fixed at the start of its Section~2. Cuntz's work on
dimension functions \cite{Cuntz1978} precedes the general criterion.
The correspondence between \(2\)-quasitraces and lower semicontinuous
dimension functions is due to Blackadar--Handelman
\cite[Theorems~II.2.2 and II.3.1]{BH}; the extension of
semigroup states underlying this formulation is discussed in
\cite{BR}. In particular, Theorem~\ref{thm:existence} has no exactness,
simplicity, or representation assumption.

For finite-rank complex vector bundles, Hermitian metrics are fixed.
Adjoints, positive endomorphisms, and their square roots are taken
fiberwise. A positive lower bound on the adjoint square of a
continuous bundle map makes it injective with a continuous image
subbundle: if \(W^*W\) is invertible, then
\[
 Q=W(W^*W)^{-1/2},\qquad Q^*Q=I,
\]
has the same image, and \(QQ^*\) is its continuous orthogonal
projection. We will use this observation when passing from
polynomial operators to the Chern-class obstruction.

\section{Two bundle-map lemmas}
\label{sec:bundle-map-lemmas}

We first construct maps whose collective images expand subspaces.  We
then normalize such a family by a fiberwise minimization.  Uniqueness of
the minimizer will make the normalization continuous on the base, even
when the bundles are nontrivial.

\begin{lemma}[Avoiding small collective images]
\label{lem:incidence}
Let $E,F$ be smooth Hermitian complex vector bundles of ranks $d,r$ over
a compact smooth manifold $Y$ of real dimension $2h$, and let $q\geq1$.
Let $\mathcal A$ be a finite set of integer pairs $(s,t)$ with
$1\leq s\leq d$ and $0\leq t\leq r$ such that
\begin{equation}
  (qs-t)(r-t)-s(d-s)>h
  \qquad ((s,t)\in\mathcal A).
  \label{eq:incidence-dimension}
\end{equation}
There are smooth bundle maps $T_1,\ldots,T_q:E\to F$ such that, for every
$y\in Y$, every $(s,t)\in\mathcal A$, and every complex $s$-plane
$S\subset E_y$,
\[
  \dim_{\mathbb C}\sum_{i=1}^q T_i(y)S>t.
\]
Here and below the sum denotes linear span.  The same statement holds
over a point, with $h=0$.
\end{lemma}

\begin{proof}
Choose a finite cover trivializing $\operatorname{Hom}(E,F)$ and a smooth
partition of unity subordinate to it, with each support contained in its
trivializing open set.  Multiply every local frame section by the
corresponding partition function and extend by zero.  Their complex
linear span is a finite-dimensional space $V$ of smooth bundle maps
whose evaluation map
\[
  V\longrightarrow\operatorname{Hom}(E_y,F_y)
\]
is surjective for every $y$: at least one partition function is nonzero
there, and its multiplied frame still spans the fiber.  Thus the space
$\mathcal V=V^q$, of some complex dimension $M$, parametrizes tuples with
surjective evaluation onto $\operatorname{Hom}(E_y,F_y)^q$.

Fix $(s,t)\in\mathcal A$.  The Grassmann bundle of triples $(y,S,U)$,
where $S\subset E_y$ is an $s$-plane and $U\subset F_y$ is a $t$-plane,
has real dimension
\[
  2h+2s(d-s)+2t(r-t).
\]
Over this bundle, evaluation followed by restriction and quotient gives
a smooth, fiberwise surjective complex-linear map from the trivial
bundle with fiber $\mathcal V$ to the bundle with fiber
$\operatorname{Hom}(S,F_y/U)^q$.  Surjectivity follows because any linear
map $S\to F_y/U$ extends to a map $E_y\to F_y$, and evaluation from $V$
is surjective.  The target has complex rank $qs(r-t)$, so the kernels
form a smooth vector bundle whose total space $\mathcal I_{s,t}$ has real
dimension
\[
  2M+2h+2s(d-s)+2t(r-t)-2qs(r-t)<2M,
\]
where the inequality is precisely \eqref{eq:incidence-dimension}.

The image of the smooth projection
$\mathcal I_{s,t}\to\mathcal V$ consists exactly of tuples having some
$s$-plane whose collective image lies in a $t$-plane.  A smooth image of a manifold of smaller dimension is a Lebesgue
null set, as in the lower-dimensional case of Sard's theorem
\cite{Sard}.  Indeed, on a compact coordinate patch the projection is
Lipschitz.  Subdivision into cubes of side $\delta$ covers its image
by $O(\delta^{-a})$ balls of radius $O(\delta)$, where $a<2M$ is
the domain dimension; their total $2M$-dimensional volume tends to
zero.  A countable exhaustion by such patches handles the possibly
noncompact total space.  The finite
union of these null sets over $\mathcal A$ cannot exhaust $\mathcal V$.
A tuple outside that union has the required properties.
\end{proof}

The next normalization uses a determinant potential. Gurvits's
operator-scaling theory uses determinant capacity to relate rank
nondecrease to approximate doubly stochastic scaling \cite[Section~4]{Gurvits2003}.
Related affine log-determinant functions are studied by Sra and
Hosseini \cite[Corollary~2.11]{SraHosseini2015}. We prove the version
needed here, including the strict convexity and continuity that allow
normalization over nontrivial bundles.

\begin{lemma}[Continuous normalization]
\label{lem:scaling}
Let $E,F$ be Hermitian complex vector bundles of positive ranks over a
compact manifold $Y$.  Suppose that continuous bundle maps
$T_1,\ldots,T_n:E\to F$ satisfy
\begin{equation}
  \dim_{\mathbb C}\sum_{i=1}^n T_i(y)S
  \geq 2\dim_{\mathbb C}S
  \qquad (y\in Y,\ S\subset E_y).
  \label{eq:expansion}
\end{equation}
Then there are continuous bundle maps $X_1,\ldots,X_n:E\to F$ with
\begin{equation}
  \sum_{i=1}^nX_i^*X_i=I_E,
  \qquad
  \sum_{i=1}^nX_iX_i^*\leq\frac23 I_F.
  \label{eq:scaled-bounds}
\end{equation}
\end{lemma}

\begin{proof}
Write $d=\operatorname{rank}E$ and $r=\operatorname{rank}F$.  In each
fiber define
\[
  D_y(P)=\sum_{i=1}^nT_i(y)P T_i(y)^*,
  \qquad
  g_y(P)=\log\det(I_{F_y}+D_y(P))-\frac32\log\det P
\]
for positive definite $P\in\operatorname{End}(E_y)$.  Both determinants
are intrinsic.  We will prove that $g_y$ has a unique minimizer depending
continuously on $y$.  The expansion factor $2$ and the coefficient
$3/2$ leave a strict margin that prevents this minimizer from escaping
to infinity.

\paragraph{Uniform compact sublevels.}
Order all eigenvalues decreasingly.  For an orthogonal rank-$s$
projection $p$ on $E_y$,
\[
  \operatorname{ran}D_y(p)=\sum_iT_i(y)\operatorname{ran}p.
\]
Indeed, both sides have orthogonal complement
$\bigcap_i\ker(pT_i(y)^*)$.  By \eqref{eq:expansion}, the $2s$-th
eigenvalue of $D_y(p)$ is positive; in particular $r\geq2d$.
Eigenvalues are continuous, and the Grassmann bundle of such projections
is compact.  Taking a minimum there and then over $1\leq s\leq d$ gives
$\epsilon>0$ such that
\[
  \mu_{2s}(D_y(p))\geq\epsilon
\]
for every $y,s,p$ under consideration.

If $P$ has eigenvalues $\lambda_1\geq\cdots\geq\lambda_d>0$, choose any
projection $p_s$ onto $s$ top eigenvectors.  Then $P\geq\lambda_s p_s$,
so positivity of $D_y$ and eigenvalue monotonicity give
\[
  \mu_{2s}(D_y(P))\geq\epsilon\lambda_s.
\]
Pairing the target eigenvalues with indices $2s-1,2s$, and discarding any
remaining nonnegative terms, yields
\begin{equation}
  \log\det(I+D_y(P))
  \geq2\sum_{s=1}^d\log(1+\epsilon\lambda_s).
  \label{eq:scaling-determinant-bound}
\end{equation}
No continuous choice of the projections $p_s$ is needed.

Set $f(t)=2\log(1+\epsilon t)-\frac32\log t$.  This function is bounded
below and tends to $+\infty$ at both endpoints of $(0,\infty)$.  Let
$b=\min f$ and $C=\max_{y\in Y}g_y(I_{E_y})$.  Equation
\eqref{eq:scaling-determinant-bound} shows that, whenever $g_y(P)\leq C$,
\[
  f(\lambda_j)\leq C-(d-1)b\qquad(1\leq j\leq d).
\]
Thus all eigenvalues of every such $P$ lie in a common compact interval
$[a,A]\subset(0,\infty)$.  Outside the compact set
$aI_{E_y}\leq P\leq AI_{E_y}$, we have
$g_y(P)>C\geq g_y(I_{E_y})$.  Minimizing on this set therefore gives a
global minimizer in each fiber, lying in the interior of the positive
definite cone.

\paragraph{Strict convexity along positive definite geodesics.}
Fix $P_0>0$ and a nonzero Hermitian $H$ in one source fiber, and put
\[
  P(u)=P_0^{1/2}\exp(uH)P_0^{1/2}.
\]
Choose an orthonormal eigenbasis $e_j$ for $H$, with eigenvalues $h_j$,
and set $v_{ij}=T_iP_0^{1/2}e_j$.  Then
\[
  D(P(u))=\sum_{i,j}e^{u h_j}v_{ij}v_{ij}^*.
\]
The Gram determinant expansion of $\det(I+CC^*)$, applied to the columns
$e^{u h_j/2}v_{ij}$, expresses $\det(I+D(P(u)))$ as a finite sum
$\sum_\alpha c_\alpha e^{u\alpha}$ with $c_\alpha\geq0$.  The empty
subset contributes $1$, and singleton subsets contribute
$\|v_{ij}\|^2e^{u h_j}$.  For every $j$, the total singleton coefficient
at $h_j$ is positive: the vector $P_0^{1/2}e_j$ is nonzero and
\eqref{eq:expansion} makes the common kernel of the maps $T_i$ trivial.
Since some $h_j\ne0$, the expansion contains positive terms at two
distinct exponents, $0$ and $h_j$.

The second derivative of the logarithm of a positive exponential sum is
the variance of its exponents under weights proportional to
$c_\alpha e^{u\alpha}$.  It is therefore strictly positive here.  On the
other hand,
\[
  \log\det P(u)=\log\det P_0+u\operatorname{tr}H
\]
is affine.  Hence $g(P(u))$ is strictly convex.  Any two distinct positive
definite matrices $P_0,P_1$ are joined by such a nonconstant curve, by
taking $H=\log(P_0^{-1/2}P_1P_0^{-1/2})$.  The minimizer is unique.

\paragraph{Continuity and gluing.}
In local unitary frames, $(y,P)\mapsto g_y(P)$ is continuous and the
minimizers lie in the common compact matrix set $aI\leq P\leq AI$.
If $y_k\to y$, every subsequential limit of their minimizers minimizes
$g_y$: pass to the limit in the minimizing inequality against any fixed
positive definite comparison matrix.  Uniqueness forces every such
limit to be the minimizer at $y$, proving continuity.  On overlaps of
unitary frames, $P$ and $I_{F_y}+D_y(P)$ change by unitary conjugation,
preserving both determinants.  Uniqueness
therefore makes the local minimizers glue to a continuous positive
definite endomorphism $P$ of $E$.

\paragraph{The normalization.}
Put $Q=I_F+D(P)$.  Differentiation at the minimizer in an arbitrary
Hermitian source direction $K$ gives
\[
  0=\operatorname{tr}\left[
     \left(\sum_iT_i^*Q^{-1}T_i-\frac32 P^{-1}\right)K
   \right].
\]
The trace pairing on Hermitian matrices is nondegenerate, so
$\sum_iT_i^*Q^{-1}T_i=\frac32P^{-1}$.  Define
\[
  X_i=\sqrt{\frac23}\,Q^{-1/2}T_iP^{1/2}.
\]
Continuous functional calculus makes these continuous bundle maps, and
the order of the factors gives
\begin{align*}
  \sum_iX_i^*X_i
  &=\frac23 P^{1/2}\left(\sum_iT_i^*Q^{-1}T_i\right)P^{1/2}
    =I_E,\\
  \sum_iX_iX_i^*
  &=\frac23 Q^{-1/2}D(P)Q^{-1/2}
    =\frac23(I_F-Q^{-1})\leq\frac23 I_F.
\end{align*}
These are \eqref{eq:scaled-bounds}.
\end{proof}

\section{Bundle tests of bounded degree}
\label{sec:bundle-tests}

We now construct the finite topological tests used to exclude polynomial
witnesses of proper infiniteness.  The matrix size $m$ and the degree bound
$R$ determine the bundles and their base, while the number of maps remains
fixed at $24$.
The use of characteristic classes to obstruct comparison follows the
methods of Villadsen \cite[Section~3]{Villadsen1999} and R{\o}rdam
\cite[Sections~3--4]{Rordam2003}. Our test retains enough of a Chern
class to exclude a specified trivial subbundle, even though the maps
between successive levels have the normalized column and row bounds.

\begin{proposition}\label{prop:bundle-test}
For every pair of integers $m,R\geq 1$, set
\begin{equation}\label{eq:test-ranks}
 d_0=64m4^R,\qquad d_j=4^j d_0\quad(0\leq j\leq R).
\end{equation}
There are a compact smooth manifold $Y$, Hermitian complex vector bundles
$E_j\to Y$ for all $j\geq0$, and continuous bundle maps
$X_i^{(j)}:E_j\to E_{j+1}$, for $1\leq i\leq24$ and $j\geq0$, such that
$E_0$ is trivial of rank $d_0$, the rank of $E_j$ is $d_j$ for
$0\leq j\leq R$, and
\begin{equation}\label{bundle:scaled-identities}
 \sum_{i=1}^{24}(X_i^{(j)})^*X_i^{(j)}=I_{E_j},
 \qquad
 \sum_{i=1}^{24}X_i^{(j)}(X_i^{(j)})^*
 \leq\frac23 I_{E_{j+1}}.
\end{equation}
Moreover,
\begin{equation}\label{eq:chern-obstruction}
 \begin{gathered}
 \left(\bigoplus_{j=0}^R E_j\right)^{\oplus m}\\
 \text{contains no trivial subbundle of rank }2md_0.
 \end{gathered}
\end{equation}
\end{proposition}

\begin{proof}
Fix the matrix size $m$ and degree bound $R$.
The bundles must satisfy two requirements.  Their sum through level $R$
must retain a Chern class after allowing the trivial directions coming
from $E_0$.  At the same time, successive bundles must admit a fixed
number of maps that expand every source subspace.  We reduce each old
line-bundle multiplicity by a factor of four at every step, so the total
trivial part stays small.  Each new projective-space factor is then sized
to retain the Chern class contributed by every later occurrence of its
line bundle.  The initial rank leaves enough room in the reduced blocks
to expand small subspaces; a large new block handles the others.

\paragraph{The bundles and their base.}
For $1\leq i\leq R$, define the fresh multiplicities
\[
 b_i=\frac{15}{4}d_{i-1}.
\]
Let the multiplicity of the line bundle $L_i$ in $E_j$ be
\begin{equation}\label{bundle:multiplicities}
 \begin{split}
 a_{0,j}&=d_0 4^{-j}=64m4^{R-j},\\
 a_{i,j}&=b_i4^{-(j-i)}
          =960m4^{R+2i-j-2}\qquad(1\leq i\leq j\leq R).
 \end{split}
\end{equation}
All these numbers are positive integers.  Put
\begin{equation}\label{eq:test-base}
 K_i=m\sum_{j=i}^R a_{i,j}
     =320m^2 4^{2i-2}\bigl(4^{R-i+1}-1\bigr),
 \qquad
 Y=\prod_{i=1}^R\mathbb{CP}^{K_i}.
\end{equation}
Thus every $K_i$ is a positive integer.  Take $L_0$ to be the trivial line
bundle on $Y$, and for $i\geq1$ let $L_i$ be the pullback of the hyperplane
line bundle from the $i$th factor.  Choose Hermitian metrics and define
\[
 E_j=\bigoplus_{i=0}^j L_i^{\oplus a_{i,j}}
 \qquad(0\leq j\leq R).
\]
In passing from $E_j$ to $E_{j+1}$, each old multiplicity is divided by
four and a fresh block of rank $b_{j+1}=(15/4)d_j$ is added.
Consequently the ranks satisfy
\[
 \operatorname{rank}E_{j+1}
 =\frac14d_j+\frac{15}{4}d_j=4d_j,
\]
which proves the rank assertion by induction.  Formula~\eqref{bundle:multiplicities}
also shows that every multiplicity $D$ to which division by four is
applied is divisible by four and satisfies $D\geq256m$.

\paragraph{The retained Chern class.}
The choice of $K_i$ records exactly the total multiplicity of $L_i$ in
the bundle to be tested.  Set
\[
 G=\left(\bigoplus_{j=0}^R E_j\right)^{\oplus m},\qquad
 h_0=m\sum_{j=0}^R a_{0,j},\qquad p=\sum_{i=1}^R K_i.
\]
The definitions of the multiplicities and $K_i$ give
\begin{equation}\label{bundle:chern-decomposition}
 G\cong L_0^{\oplus h_0}\oplus
             \bigoplus_{i=1}^R L_i^{\oplus K_i},\qquad
 h_0=\frac43md_0\bigl(1-4^{-(R+1)}\bigr)<2md_0.
\end{equation}
Let $z_i=c_1(L_i)$.  The integral cohomology ring is
\[
 H^*(Y;\mathbb Z)
 \cong\mathbb Z[z_1,\ldots,z_R]/
       (z_1^{K_1+1},\ldots,z_R^{K_R+1}),\qquad \deg z_i=2.
\]
In particular $\prod_i z_i^{K_i}$ is nonzero.  The Whitney sum formula
and the normalization for line bundles give
\[
 c_p(G)=\prod_{i=1}^R z_i^{K_i}\ne0;
\]
we use the standard Chern-class properties in
\cite[Theorem~3.2]{HatcherVB}.
If $G$ had a trivial subbundle of rank $k>h_0$, its Hermitian complement
would have rank $h_0+p-k<p$.  The Whitney sum formula would identify
$c_p(G)$ with the $p$th Chern class of that complement, which vanishes
above its rank.  This contradiction, together with
\eqref{bundle:chern-decomposition}, proves \eqref{eq:chern-obstruction}.
The argument applies to any continuous subbundle, without a condition
on its position relative to the displayed line summands.

\paragraph{The bases for the expanding maps.}
It remains to construct the maps between these bundles.  For the step
$j\to j+1$, with $0\leq j<R$, the source and target use only
$L_0,\ldots,L_{j+1}$ and are pulled back from
\[
 Y_{j+1}=\prod_{i=1}^{j+1}\mathbb{CP}^{K_i}.
\]
The dimensions $K_i$ here are those in \eqref{eq:test-base}, including
all later occurrences of $L_i$ through level $R$.  The complex dimension
$\eta_j$ of this smaller base satisfies
\begin{equation}\label{bundle:prefix-dimension}
 \begin{split}
 \eta_j=\sum_{i=1}^{j+1}K_i
 &\leq\frac43m\sum_{i=1}^{j+1}b_i\\
 &=\frac53md_0(4^{j+1}-1)
 <\frac{20}{3}md_j.
 \end{split}
\end{equation}
Here the first inequality uses the geometric series with ratio $1/4$.
We construct the maps for this step on $Y_{j+1}$, and then pull them back
to $Y$.  The estimate by a multiple of $md_j$ is what will allow the
incidence argument to handle all subspaces of dimension greater than
$m$, even though the full base $Y$ may be much larger.

\paragraph{Sixteen maps for small subspaces.}
Decompose $E_{j+1}$ as the sum of its replacement blocks and its fresh
block.  On a source block $L_i^{\oplus D}$, the corresponding replacement
block is $L_i^{\oplus(D/4)}$.  We seek $16$ scalar matrices between the
multiplicity spaces and tensor them with the identity of $L_i$.
For $1\leq s\leq m$, set $r=D/4$ and $t=2s$.  Since $D\geq64m$,
\[
 t<r,\qquad r-2s\geq D/8.
\]
The incidence expression of Lemma~\ref{lem:incidence}, with $q=16$ and
base a point, is strictly positive:
\begin{equation}\label{bundle:small-incidence}
 (16s-2s)(D/4-2s)-s(D-s)
 =\frac52sD-27s^2
 \geq133ms>0.
\end{equation}
Applying Lemma~\ref{lem:incidence} simultaneously for $s=1,\ldots,m$ gives matrices
whose images of every such $s$-plane span a space of dimension at least
$2s$.  In fact, avoiding dimension at most $2s$ gives a slightly stronger
conclusion.  Make these choices for each block and assemble them into
$16$ block diagonal maps $T_1,\ldots,T_{16}$ to the replacement part,
with zero component in the fresh block.

We verify that the same $16$ maps expand every subspace of the whole
source of dimension at most $m$.  Work in one fiber and write the
ordered source and replacement blocks as $V_k$ and $F_k$.  Let
$V_{\leq k}$ and $F_{\leq k}$ denote their initial sums.  Given a subspace
$S$ of dimension at most $m$, set
\[
 S_k=S\cap V_{\leq k},\qquad
 W=\sum_{i=1}^{16}T_iS,\qquad
 W_k=W\cap F_{\leq k}.
\]
Use the zero spaces before the first flag terms.  Projection onto the
$k$th block identifies $S_k/S_{k-1}$ and $W_k/W_{k-1}$ with subspaces
of $V_k$ and $F_k$, respectively.  Block diagonality gives
$T_i(S_k)\subseteq W_k$ and hence
\[
 \sum_{i=1}^{16}T_{i,k}(S_k/S_{k-1})
 \subseteq W_k/W_{k-1},
\]
where $T_{i,k}$ is the map on the $k$th block.  Each source quotient has
dimension at most $m$, so the block estimates yield
\begin{equation}\label{bundle:flag-expansion}
 \dim W
 =\sum_k\dim(W_k/W_{k-1})
 \geq2\sum_k\dim(S_k/S_{k-1})
 =2\dim S.
\end{equation}
This proves the required estimate even when $S$ has correlated
components in several blocks.

\paragraph{Eight maps for large subspaces.}
Let $d=d_j$ and let the target now be the fresh block of rank
$r=(15/4)d$.  For every integer $m<s\leq d$, take $q=8$ and $t=2s$.
The rank endpoint is valid because $2s\leq2d<r$.  The incidence
expression satisfies
\begin{equation}\label{bundle:large-incidence}
 \begin{split}
 (8s-2s)(15d/4-2s)-s(d-s)
 &=\frac{43}{2}sd-11s^2\\
 &\geq\frac{21}{2}sd
 >\frac{20}{3}md
 >\eta_j.
 \end{split}
\end{equation}
The first inequality follows from $s\leq d$.  Lemma~\ref{lem:incidence} on $Y_{j+1}$, applied to the finite
list $s=m+1,\ldots,d$, therefore supplies eight smooth maps to the fresh
block whose image spans have dimension at least $2s$ for all these
subspaces.  Regard these as maps to $E_{j+1}$ with zero replacement
components.

Together, the two lists give $24$ maps $T_i:E_j\to E_{j+1}$ satisfying
\[
 \dim\sum_{i=1}^{24}T_iS\geq2\dim S
\]
for every fiber subspace $S$.  The estimate is supplied by the first
list when $\dim S\leq m$ and by the second when $\dim S>m$; adding maps
cannot decrease an image span.  Pull back to $Y$ and apply
Lemma~\ref{lem:scaling}.  This gives continuous maps $X_i^{(j)}$
satisfying \eqref{bundle:scaled-identities} for every $0\leq j<R$.
The scaling may mix the displayed line summands, but each resulting
map still has source $E_j$ and target $E_{j+1}$.

\paragraph{Continuation beyond the test.}
For $j\geq R$, define $E_{j+1}=E_j\oplus E_j$, and let
\[
 X_1^{(j)}=2^{-1/2}\iota_1,\qquad
 X_2^{(j)}=2^{-1/2}\iota_2,\qquad
 X_i^{(j)}=0\quad(3\leq i\leq24),
\]
where $\iota_1,\iota_2$ are the two inclusions.  Their column sum is
$I_{E_j}$ and their row sum is $\tfrac12I_{E_{j+1}}$, so
\eqref{bundle:scaled-identities} holds for these steps as well.
This continuation leaves the retained bundle in
\eqref{eq:chern-obstruction} unchanged and completes the construction.
\end{proof}

\section{The algebra and a uniform additivity defect}
\label{sec:algebra}

We now put all the bundle tests into one algebra.  The uniform number of
maps and their uniform norm bounds allow this construction, while each
test obstructs polynomial candidates at its prescribed matrix size and
degree.

\begin{theorem}
\label{thm:construction}
There is a separable unital complex $C^*$-algebra
$B=C^*(1,x_1,\ldots,x_{24})$ such that
\begin{equation}
  \sum_{i=1}^{24}x_i^*x_i=1,
  \qquad
  \sum_{i=1}^{24}x_ix_i^*\le\frac23\,1,
  \label{eq:algebra-relations}
\end{equation}
and $M_m(B)$ is not properly infinite for every integer $m\ge1$.
\end{theorem}

\begin{proof}
For each pair of integers $m,R\ge1$, choose the data supplied by
Proposition~\ref{prop:bundle-test}, and denote its base by $Y_{m,R}$.
At a point $y\in Y_{m,R}$ form the Hilbert space
\[
  \mathcal H_{m,R,y}=\bigoplus_{j=0}^{\infty}E_j(y).
\]
For $1\leq i\leq24$, let $S_i^{m,R}(y)$ be the forward shift on this
space whose block from $E_j(y)$ to $E_{j+1}(y)$ is $X_i^{(j)}(y)$.
The column identity in
Proposition~\ref{prop:bundle-test} gives
\[
  X_i^{(j)*}X_i^{(j)}\le I_{E_j},
\]
so every block, and hence every shift, has norm at most one.  These
bounds hold uniformly over all coordinates.

Let $\mathcal D$ be the bounded $C^*$-product of the algebras
$\mathcal B(\mathcal H_{m,R,y})$ over all these indices: its elements are
the operator families of finite supremum norm.  The uniform bound
therefore defines a single tuple in $\mathcal D$,
\[
  x_i=\bigl(S_i^{m,R}(y)\bigr)_{m,R\geq1,\;y\in Y_{m,R}}
  \qquad(1\leq i\leq24).
\]
The column sum of the shifts is the identity on every level.  Their
row sum is zero on level zero and, on level $j\ge1$, is
\[
  \sum_{i=1}^{24}X_i^{(j-1)}X_i^{(j-1)*}
  \le\frac23 I_{E_j}.
\]
Thus \eqref{eq:algebra-relations} holds in $\mathcal D$.  Take
$B=C^*(1,x_1,\ldots,x_{24})\subseteq\mathcal D$.  This is a nonzero
unital algebra, and it is separable: the $*$-polynomials in these
generators with rational complex coefficients form a countable dense
subset.

Fix $m\ge1$ and suppose, towards a contradiction, that $M_m(B)$
contains isometries $s_1,s_2$ with orthogonal ranges.  Their rectangular
concatenation satisfies
\[
  W=[s_1\ s_2]\in M_{m,2m}(B),
  \qquad W^*W=I_{2m}.
\]
By norm density choose $P\in M_{m,2m}(B)$ with $*$-polynomial entries
such that $\|P-W\|<1/5$.  Since $\|W\|=1$,
\[
  \|P^*P-I_{2m}\|
  \le 2\|P-W\|+\|P-W\|^2<\frac12,
\]
and consequently
\begin{equation}
  P^*P\ge\frac12 I_{2m}.
  \label{eq:polynomial-lower-bound}
\end{equation}
After fixing this approximation, choose $R\ge1$ at least the length of
every word occurring in its entries.  We use the coordinate family
belonging to the already constructed test $(m,R)$.

A word of length at most $R$, applied to level zero, takes values in
levels zero through $R$.  Each generator raises the level by one, each
adjoint lowers it by one, and an adjoint applied at level zero gives
zero.  Constant terms preserve level zero.  Restricting the evaluated
matrix $P$ to the input at level zero therefore gives a fiberwise map
\begin{equation}
  P_0:E_0^{\oplus2m}\longrightarrow
  F_{m,R}:=\left(\bigoplus_{j=0}^{R}E_j\right)^{\oplus m}
  \quad\text{over }Y_{m,R}.
  \label{eq:polynomial-bundle-map}
\end{equation}
This is a continuous bundle map.  Indeed, each block of the
restriction is a finite sum of compositions of the continuous maps
$X_i^{(j)}$ and their adjoints, with $j<R$.  This description is
intrinsic to the bundles; no simultaneous trivialization of the
Hilbert fibers or continuity between different tests is required.
The continuation beyond level $R$ does not enter these compositions. Figure~\ref{fig:finite-test} illustrates
this finite propagation and the resulting bundle map.

\begin{figure}[htbp]
\centering
\begin{tikzpicture}[>=Latex, font=\small,
  level/.style={draw, rounded corners=2pt, minimum width=12mm,
                 minimum height=9mm, fill=blue!4}]
  \node[level] (e0) at (0,0) {$E_0$};
  \node[level] (e1) at (1.8,0) {$E_1$};
  \node (dots) at (3.45,0) {$\cdots$};
  \node[level] (er) at (5.1,0) {$E_R$};
  \node[level,fill=gray!8] (beyond) at (7.2,0) {$E_{R+1}$};
  \draw[->] (e0.north east) to[bend left=15] node[above] {$x_i$} (e1.north west);
  \draw[->] (e1.south west) to[bend left=15] node[below] {$x_i^*$} (e0.south east);
  \draw[<->] (e1) -- (dots);
  \draw[<->] (dots) -- (er);
  \draw[<->,dashed] (er) -- (beyond);
  \begin{scope}[on background layer]
    \node[draw=blue!50!black,dashed,rounded corners=3pt,
          inner xsep=4mm,inner ysep=8mm,fit=(e0)(er)] (window) {};
  \end{scope}
  \node[anchor=south,align=center] at (2.55,1.5)
    {A word of length at most $R$, starting at level zero,\\ stays in these levels};
  \node[anchor=north,align=center] at (2.55,-1.5)
    {$P_0:E_0^{\oplus 2m}\longrightarrow
       \bigl(E_0\oplus\cdots\oplus E_R\bigr)^{\oplus m}$};
\end{tikzpicture}
\caption{The test indexed by $(m,R)$, shown schematically. A polynomial
matrix $P$ whose entries contain only words of length at most $R$ maps
its trivial level-zero input into the displayed target. If
$P^*P\ge\tfrac12 I$, this restriction would embed a trivial bundle of
rank $2md_0$ into that target,
contradicting its retained Chern class. The arrows show level changes;
each forward block is one of the maps $X_i^{(j)}$. The continuation after
level $R$ does not affect this test. When $R=1$, the labels $E_1$ and
$E_R$ denote the same level.}
\label{fig:finite-test}
\end{figure}


At each $y$, let $\iota_0$ include the level-zero input in
$\mathcal H_{m,R,y}^{\oplus2m}$, and let $\iota_F$ include
$F_{m,R}(y)$ in $\mathcal H_{m,R,y}^{\oplus m}$.  The range statement
above says $P(y)\iota_0=\iota_F P_0(y)$.  Evaluation preserves
\eqref{eq:polynomial-lower-bound}, so compression of this quadratic-form
inequality gives
\[
  P_0(y)^*P_0(y)
  =\iota_0^*P(y)^*P(y)\iota_0
  \geq\frac12 I
  \qquad(y\in Y_{m,R}).
\]
Hence
\[
  Q=P_0(P_0^*P_0)^{-1/2}
\]
is a continuous fiberwise isometry.  Its image is a subbundle of
$F_{m,R}$ isomorphic to the trivial bundle $E_0^{\oplus2m}$, of rank
$2m\operatorname{rank}(E_0)$.  This contradicts
Proposition~\ref{prop:bundle-test}.  Since $m$ was arbitrary, the
conclusion holds for every matrix algebra over $B$.
\end{proof}

The passage from the two relations to a specified positive pair is
useful independently of the construction.  It requires only the
$1$-quasitrace axioms once a quasitrace on the matrix algebra is given.

\Needspace{12\baselineskip}
\begin{lemma}
\label{lem:defect}
Let $B$ be a unital complex $C^*$-algebra, let $N\ge1$, and suppose
$x_1,\ldots,x_N\in B$ satisfy
\[
  \sum_{i=1}^N x_i^*x_i=1,
  \qquad
  \sum_{i=1}^N x_ix_i^*\le c1
  \quad\text{for some }0\le c<1.
\]
There are positive contractions $a,b\in M_{2N}(B)$, determined by this
tuple, such that every normalized $1$-quasitrace $\tau$ on $M_{2N}(B)$
satisfies
\begin{equation}
  \tau(a+b)-\tau(a)-\tau(b)\ge\frac{1-c}{2N}.
  \label{eq:uniform-defect}
\end{equation}
\end{lemma}

\begin{proof}
Put
\[
  u_i=x_i^*x_i,\qquad v_i=x_ix_i^*,\qquad
  H_i=\sum_{j=1}^i u_j,\qquad J_i=\sum_{j=1}^i v_j,
\]
with $H_0=J_0=0$.  We use these partial sums to define two diagonal
matrices.  In their additivity defect, the first $N$ blocks will telescope
the column partial sums, while the last $N$ use complements to telescope
the row partial sums with the opposite sign and keep the entries positive.
\begin{equation}
\begin{aligned}
  a&=\operatorname{diag}
      (H_0,\ldots,H_{N-1},\,1-J_0,\ldots,1-J_{N-1}),\\
  b&=\operatorname{diag}
      (u_1,\ldots,u_N,\,1-v_1,\ldots,1-v_N).
\end{aligned}
\label{eq:defect-pair}
\end{equation}
All entries are positive contractions: $0\le H_i\le H_N=1$,
$0\le u_i\le1$, and $0\le J_i,v_i\le J_N\le c1$.

Fix a normalized $1$-quasitrace $\tau$ on $M_{2N}(B)$.  Write
$\iota_k(z)$ for $z\in B$ placed in the $k$th diagonal corner and
define $\varphi(z)=\tau(\iota_1(z))$.  If $z\ge0$, let $t$ be the
matrix with $z^{1/2}$ in its $(k,\ell)$ entry and zeros elsewhere.
Then $t^*t=\iota_\ell(z)$ and $tt^*=\iota_k(z)$, so the identity
$\tau(t^*t)=\tau(tt^*)$ gives
\[
  \tau(\iota_k(z))=\varphi(z)
  \qquad (z\ge0, 1\le k\le2N).
\]
Since self-adjoint elements supported in distinct diagonal corners
commute, abelian linearity yields
\begin{equation}
  \tau\bigl(\operatorname{diag}(z_1,\ldots,z_{2N})\bigr)
    =\sum_{k=1}^{2N}\varphi(z_k)
  \quad (z_k\ge0),
  \qquad \varphi(1)=\frac1{2N}.
  \label{eq:corner-normalization}
\end{equation}
The last equality follows by taking every $z_k=1$.  The square axiom
applied within the first corner also gives
\begin{equation}
  \varphi(u_i)=\varphi(v_i)\qquad (1\le i\le N).
  \label{eq:corner-square-equality}
\end{equation}

The first $N$ diagonal blocks contribute to the additivity defect the
terms
\[
  \varphi(H_i)-\varphi(H_{i-1})-\varphi(u_i).
\]
The last $N$ contribute
\begin{align*}
  &\varphi(2-J_i)-\varphi(1-J_{i-1})-\varphi(1-v_i)\\
  &\hspace{2em}=-\varphi(J_i)+\varphi(J_{i-1})+\varphi(v_i).
\end{align*}
Here each scalar-translation identity is taken within the abelian
algebra generated by the corner identity and the indicated
self-adjoint element.  Thus these equalities use no additivity on
noncommuting elements of $B$.  Summing the scalar defects,
telescoping, and using \eqref{eq:corner-square-equality}, we obtain
\begin{align*}
  \tau(a+b)-\tau(a)-\tau(b)
    &=\varphi(H_N)-\sum_{i=1}^N\varphi(u_i)
      -\varphi(J_N)+\sum_{i=1}^N\varphi(v_i)\\
    &=\varphi(1)-\varphi(J_N).
\end{align*}
The restriction of $\varphi$ to $C^*(1,J_N)\subseteq B$ is positive
and linear.  Since $0\le J_N\le c1$, it follows that
$\varphi(J_N)\le c\varphi(1)$.  Equation~\eqref{eq:corner-normalization}
now gives \eqref{eq:uniform-defect}.
\end{proof}

\Needspace{8\baselineskip}
\begin{proof}[Proof of Theorem~\ref{thm:main}]
Theorem~\ref{thm:construction} supplies the separable unital algebra
\(B\), generated by twenty-four elements satisfying
\eqref{eq:algebra-relations}, with no properly infinite matrix
algebra over it. Consequently the same is true of \(A=M_{48}(B)\),
because \(M_k(A)\cong M_{48k}(B)\) for every \(k\).
Theorem~\ref{thm:existence} gives a normalized \(2\)-quasitrace on \(A\).

Apply Lemma~\ref{lem:defect} with \(N=24\) and \(c=2/3\).
It supplies positive contractions \(a,b\in A\) for which every
normalized \(2\)-quasitrace has defect at least
\((1-c)/(2N)=1/144\). Matrix amplification preserves separability
and unitality, so all assertions of Theorem~\ref{thm:main} follow.
\end{proof}

\section{Tensor products and stable finiteness}
\label{sec:consequences}

The consequences in Corollary~\ref{cor:stable-finiteness} use the
established factor form of the quasitrace problem and a trace-obstruction
theorem for tensor products. We record their scopes before applying them
to the different algebras in the corollary.

\paragraph{The factor input.}
The assertion that every normalized \(2\)-quasitrace on every unital
\(C^*\)-algebra is a tracial state is equivalent to the assertion that
every \(AW^*\)-factor of type \(\mathrm{II}_1\) is a \(W^*\)-algebra;
see \cite[Section~2]{MR} and \cite[Theorem~4.11]{Gow}.
Every such factor admits
a normalized \(2\)-quasitrace, and it admits a tracial state exactly
when it is a \(W^*\)-algebra \cite[Section~2]{MR}. A finite
\(AW^*\)-factor is simple as a \(C^*\)-algebra
\cite[Remark~2.7]{Gow}. For a simple unital \(C^*\)-algebra, absence of
stable proper infiniteness is equivalent to stable finiteness
\cite[Section~2]{MR}. These statements impose no separability condition
on the factor.

\paragraph{The tensor input and the passage to \(\mathbb F_2\).}
Haagerup's trace-obstruction theorem, in the properly infinite
formulation recorded by Milh{\o}j and R{\o}rdam
\cite[Theorem~3.1]{MR}, says that a unital \(C^*\)-algebra \(D\) has no
tracial state if and only if
\(D\otimes_{\min}C_r^*(\mathbb F_\infty)\) is properly infinite.
Haagerup notes that the subgroup inclusion
\(\mathbb F_\infty\hookrightarrow\mathbb F_2\) induces a unital
inclusion \(C_r^*(\mathbb F_\infty)\hookrightarrow C_r^*(\mathbb F_2)\)
\cite[Remark~2.5(a)]{Haagerup}. The minimal tensor product preserves
injective \(*\)-homomorphisms, so it gives a unital inclusion
\[
D\otimes_{\min}C_r^*(\mathbb F_\infty)
\hookrightarrow D\otimes_{\min}C_r^*(\mathbb F_2).
\]
The two isometries witnessing proper infiniteness retain their relations
under this inclusion. Consequently,
\begin{equation}
 D\text{ has no tracial state}
 \quad\Longrightarrow\quad
 D\otimes_{\min}C_r^*(\mathbb F_2)\text{ is properly infinite}.
 \label{eq:traceless-tensor}
\end{equation}

\begin{proof}[Proof of Corollary~\ref{cor:stable-finiteness}]
For (i), Theorem~\ref{thm:main} and the factor input give an
\(AW^*\)-factor \(M\) of type \(\mathrm{II}_1\) that is not a
\(W^*\)-algebra. Hence \(M\) has no tracial state. It is simple and
admits a normalized \(2\)-quasitrace. Theorem~\ref{thm:existence}
therefore excludes stable proper infiniteness, and the criterion for
simple algebras makes \(M\) stably finite.

Powers's theorem \cite{Powers} gives simplicity of
\(R=C_r^*(\mathbb F_2)\). Its canonical tracial state is faithful, as
are the matrix traces induced by it. A faithful trace rules out a
nonunitary isometry, so \(R\) is stably finite. It is also unital.
Applying \eqref{eq:traceless-tensor} to \(M\) proves (i):
\(M\otimes_{\min}R\) is properly infinite.

For (ii), take the explicit algebra \(B\) of
Theorem~\ref{thm:construction}. No matrix algebra over \(B\) is properly
infinite, so Theorem~\ref{thm:existence} supplies a normalized
\(2\)-quasitrace. The relations \eqref{eq:algebra-relations} exclude a
tracial state: applying one to their two sums would give \(1\le2/3\).
Equation~\eqref{eq:traceless-tensor} makes
\(B\otimes_{\min}R\) properly infinite, and
Theorem~\ref{thm:existence} then excludes a normalized \(2\)-quasitrace
on that tensor product.

For (iii), Haagerup's finite trace-obstruction criterion
\cite[Lemma~2.1]{Haagerup} supplies \(y_1,\ldots,y_n\in M\) such that
\[
 \sum_{i=1}^n y_i^*y_i=1,
 \qquad
 \left\|\sum_{i=1}^n y_iy_i^*\right\|<1.
\]
Let \(C=C^*(1,y_1,\ldots,y_n)\subset M\). This algebra is separable and
unital. It inherits stable finiteness from \(M\): an isometry in
\(M_k(C)\) is an isometry in \(M_k(M)\), hence is unitary there and
therefore already unitary in \(M_k(C)\). A tracial state \(\tau\) on
\(C\) would satisfy
\[
 1=\tau\left(\sum_{i=1}^n y_i^*y_i\right)
  =\tau\left(\sum_{i=1}^n y_iy_i^*\right)
  \le\left\|\sum_{i=1}^n y_iy_i^*\right\|<1,
\]
a contradiction. This proves (iii).
\end{proof}

\begin{thebibliography}{99}

\bibitem{BH}
Bruce Blackadar and David Handelman,
\emph{Dimension functions and traces on $C^*$-algebras},
J. Funct. Anal. \textbf{45} (1982), no.~3, 297--340.
\href{https://doi.org/10.1016/0022-1236(82)90009-X}
{doi:10.1016/0022-1236(82)90009-X}.

\bibitem{BR}
Bruce Blackadar and Mikael R{\o}rdam,
\emph{Extending states on preordered semigroups and the existence of
quasitraces on $C^*$-algebras},
J. Algebra \textbf{152} (1992), no.~1, 240--247.
\href{https://doi.org/10.1016/0021-8693(92)90098-7}
{doi:10.1016/0021-8693(92)90098-7}.

\bibitem{Cuntz1978}
Joachim Cuntz,
\emph{Dimension functions on simple $C^*$-algebras},
Math. Ann. \textbf{233} (1978), no.~2, 145--153.
\href{https://doi.org/10.1007/BF01421922}{doi:10.1007/BF01421922}.

\bibitem{Gow}
Alec Gow,
\emph{On the quasitrace problem and a characterization of $W^*$-algebras},
preprint, 2026, version~2, 14~July~2026.
\href{https://arxiv.org/abs/2601.04431v2}{arXiv:2601.04431v2}.

\bibitem{Gurvits2003}
Leonid Gurvits,
\emph{Classical deterministic complexity of Edmonds' problem and quantum entanglement},
preprint, 2003.
\href{https://arxiv.org/abs/quant-ph/0303055v1}{arXiv:quant-ph/0303055v1}.

\bibitem{Haagerup}
Uffe Haagerup,
\emph{Quasitraces on exact $C^*$-algebras are traces},
C. R. Math. Rep. Acad. Sci. Canada \textbf{36} (2014), no.~2--3, 67--92.
\href{https://arxiv.org/abs/1403.7653}{arXiv:1403.7653}.

\bibitem{HatcherVB}
Allen Hatcher,
\emph{Vector Bundles and $K$-Theory},
online manuscript, version~2.2, November~2017.
\href{https://pi.math.cornell.edu/~hatcher/VBKT/VB.pdf}
{Author's manuscript}.

\bibitem{MR}
Henning Olai Milh{\o}j and Mikael R{\o}rdam,
\emph{Around traces and quasitraces},
preprint, 2023, version~2, 30~November~2023;
to appear in M\"unster J. Math.
\href{https://arxiv.org/abs/2309.17412v2}{arXiv:2309.17412v2}.

\bibitem{Powers}
Robert T. Powers,
\emph{Simplicity of the $C^*$-algebra associated with the free group on
two generators},
Duke Math. J. \textbf{42} (1975), no.~1, 151--156.
\href{https://doi.org/10.1215/S0012-7094-75-04213-1}
{\mbox{doi:10.1215/S0012-7094-75-04213-1}}.

\bibitem{Rordam2003}
Mikael R{\o}rdam,
\emph{A simple $C^*$-algebra with a finite and an infinite projection},
Acta Math. \textbf{191} (2003), no.~1, 109--142.
\href{https://doi.org/10.1007/BF02392697}{doi:10.1007/BF02392697}.

\bibitem{Sard}
Arthur Sard,
\emph{The measure of the critical values of differentiable maps},
Bull. Amer. Math. Soc. \textbf{48} (1942), no.~12, 883--890.
\href{https://doi.org/10.1090/S0002-9904-1942-07811-6}
{doi:10.1090/S0002-9904-1942-07811-6}.

\bibitem{SraHosseini2015}
Suvrit Sra and Reshad Hosseini,
\emph{Conic geometric optimization on the manifold of positive definite matrices},
SIAM J. Optim. \textbf{25} (2015), no.~1, 713--739.
\href{https://doi.org/10.1137/140978168}{doi:10.1137/140978168}.

\bibitem{Villadsen1999}
Jesper Villadsen,
\emph{On the stable rank of simple $C^*$-algebras},
J. Amer. Math. Soc. \textbf{12} (1999), no.~4, 1091--1102.
\href{https://doi.org/10.1090/S0894-0347-99-00314-8}
{doi:10.1090/S0894-0347-99-00314-8}.

\end{thebibliography}

\end{document}
